\chapter*{Manifeste : Le Chaos et la Jubilation}
\addcontentsline{toc}{chapter}{Manifeste : Le Chaos et la Jubilation}

\textit{Nous n’avons pas créé une maison d’édition. Nous avons lancé un mouvement.}

\vspace{1cm}

Nous refusons le conformisme intellectuel. Cette habitude molle qui consiste à répéter des opinions prémâchées en les confondant avec des convictions.

Nous refusons la résignation. Le fatalisme déguisé en lucidité. Le « c’est trop tard » qui évite d’agir.

Nous refusons le monopole du logos. L’idée selon laquelle seule l’analyse rationnelle produirait du savoir légitime. Nous ne combattons pas la raison. Nous combattons son hégémonie.

Car une vérité demeure, qui est trop souvent tue : la pensée créative a la même dignité que la pensée logique. C'est la fabrique de l'avenir. Et fabricant l'avenir, elle crée les humains qui vont y vivre. Elle donne des armes intellectuelles pour changer l'époque.

\section*{La Jubilation}
Nous voulons provoquer une expérience jubilatoire. La jubilation, c’est le frisson de comprendre quelque chose que nous savions sans qu'on nous l'ait jamais dit. C'est l'émotion de saisir une chose complexe qui nous échappait. C'est la "libido sciendi" : la joie intellectuelle qui circule dans notre corps.

Un livre \textbf{NIV} doit rendre son lecteur plus libre, plus joyeux, plus vivant.

\section*{Embrasser le Chaos}
Beaucoup cherchent à le dompter. D’autres préfèrent le fuir. Nous proposons autre chose : l’embrasser. L’aimer. Le rendre fécond. Nos livres sont des substances actives. Des rites de passage. Des outils pour habiter l’époque sans céder au cynisme.

\section*{L'Impossible}
Notre ambition est spirituelle. Nous voulons que chacun choisisse son impossible : un impossible qui le touche et le passionne. Nous voulons qu'il s'attaque à cet impossible sans pitié. Et puis, lorsqu'il aura vaincu, nous voulons qu'il nous raconte cette histoire.

\vspace{1.5cm}

\begin{center}
\textit{Alors, de quel côté de la force êtes-vous ? \\
Du côté de la répétition qui rassure ? Ou du côté de la pensée qui jubile ?}
\end{center}
